

% DOCUMENTO OSM - NA CONTEXTUALIZAÇAO PROMOÇAO A SAUDE MENTAL
% FALAR DA VIGILANCIA NA CONTEXTUALIZAÇÃO
%desatualizado fonaprace é precario usar outra referencia


% De acordo com o Fórum Nacional de Pró-reitores de Assuntos Comunitários e Estudantis (FONAPRACE) em 2019, o impacto da solidão e depressão e percepção ruim do curso e relacionamento com colegas de sala aumentam o risco de desistência, reprovações e trancamentos na universidade.

Devido aos altos índices de risco psicológico destacados por \citeonline{silveira_saumental_2011}, o ambiente universitário revela-se potencialmente hostil à saúde mental dos estudantes, exigindo a implementação de estratégias de promoção à saúde que vão além de ações institucionais, como palestras, e incluam também abordagens personalizadas e voltadas ao indivíduo, como diários de humor \cite{schueller_2019}.

Nesse contexto, a vigilância participativa configura-se como uma abordagem essencial no monitoramento da população acadêmica, ao permitir a identificação precoce e a resposta eficiente a situações que afetam coletivamente a saúde mental \cite{lealneto2016digital}. Sendo assim, a fim de fornecer uma maneira de auxiliar a promoção de saúde nas IES, é fundamental a utilização deste conceito.

\vspace{23em}

Além disso, a adoção de hábitos saudáveis mostra-se eficaz no enfrentamento e na mitigação dos impactos das doenças mentais, conforme apontado por \citeonline{barroso_solidao_2019}. Ou seja, A promoção desses hábitos pode refletir positivamente no desempenho acadêmico e na permanência dos estudantes nos cursos. Além disso, com base no conceito de UPS se da a necessidade de criar meios para auxiliar a promoção a saúde nas IES, como mencionado por \citeonline{rezende_promocao_2022}, não há uma fórmula única para promover a saúde no ambiente universitário.

Dessa forma, é fundamental a criação de uma ferramenta para auxiliar a adoção de hábitos saúdaveis e promover a saúde mental no meio acadêmico e uma ferramenta para auxiliar profissionais de saúde acadêmicos a identificar alunos em risco. Ademais, tendo em vista o conceito de suporte social é possível a utilização da tecnologia para criação de um software que realize esse papel \cite{siqueira_construcao_2008}. 



% Sendo assim, com base no conceito de ups se da a necessidade de criar meios para auxiliar a ps na faculdade
% Ademais segundo... Se da necessidade da criaçao de uma aplicaçao para auxiliar a triagem de alunos